\chapter{Deploying Post-Quantum Cryptography}
\label{ch:applications}

\section{Why this chapter matters}
The previous parts built and analysed the algorithms. This chapter is about the harder half of the problem: getting them into real systems -- web protocols, digital certificates, and even smart contracts -- correctly and in time. The forcing function is \emph{harvest now, decrypt later} (Chapter~\ref{ch:quantum}): an adversary can record encrypted traffic today and decrypt it once a quantum computer exists, so any long-lived secret is already exposed. Migration must therefore begin before the threat arrives, not after, and migration is an engineering project at least as much as a mathematical one.

\section{Learning goals}
After reading this chapter, you should be able to:

\begin{itemize}
  \item explain why real deployments use \emph{hybrid} classical-plus-postquantum cryptography, what crypto-agility means, and when hybrid is the wrong choice;
  \item describe how TLS 1.3 carries a hybrid key exchange and post-quantum certificates, and what the size costs are;
  \item say what has already shipped, and why key exchange migrated years ahead of signatures;
  \item explain why verifying post-quantum signatures on-chain is expensive and how projects reduce that cost;
  \item weigh the size and performance trade-offs that drive scheme choice in practice.
\end{itemize}

\section{Hybrid cryptography and crypto-agility}
The safest way to deploy a young algorithm is not to trust it alone. A \emph{hybrid} scheme combines a well-understood classical algorithm with a post-quantum one through an analysed protocol construction.

\begin{definition}[Hybrid key exchange]
A hybrid key exchange combines a classical key-agreement contribution (for example
X25519 ECDHE) with a post-quantum KEM contribution (for example ML-KEM-768) through
an analysed protocol construction. X25519 is Diffie--Hellman key agreement, not a
KEM. A robust construction binds algorithm identifiers, peer identities, and the
handshake transcript in a KDF extract/expand step. Its design goal is to retain
session-key security when at least one component remains secure, subject to the
combiner and protocol's required conditions; arbitrary hash concatenation is not a
general theorem.
\end{definition}

Hybrid deployment hedges two opposite risks at once: a future quantum computer breaking the classical part, and an as-yet-undiscovered classical weakness in the still-young post-quantum part. It is the mainstream migration strategy precisely because it removes the need to bet everything on one new assumption.

\begin{remark}[Crypto-agility]
Hybrids only work if systems can \emph{negotiate} which algorithms they use rather than hardcoding one. \emph{Crypto-agility} -- designing protocols, certificate formats, and code so the algorithm can be swapped without re-architecting -- is the quiet prerequisite for the whole migration. The schemes will keep evolving; the plumbing must not assume any single one.
\end{remark}

\subsection{Where hybrid does not fit}
\label{sec:hybrid-limits}

Hybrid is the right default for the protocols above, and it is worth being
explicit about \emph{why}, because the reasons are properties of those protocols
rather than facts about cryptography --- and one important class of system has
none of them.

In TLS, a hybrid transition is nearly free. The extra classical key exchange
costs a few hundred bytes on a connection that is discarded seconds later, both
endpoints negotiate the algorithm without asking anyone's permission, and no user
is aware the change happened. The cost is paid once, by software vendors, and it
is invisible.

On a blockchain, every one of those properties fails. The extra signature is not
transient: it is written permanently into a block, relayed to every node, and
stored by every archive, so a hybrid transaction costs its wire size
\emph{forever} and to \emph{everyone}. There is no negotiation, because consensus
rules are the same for all participants. And the change is emphatically not
invisible --- users must act, first to move to the hybrid scheme and then again
to move off it, which means the ecosystem pays the coordination cost of a
migration \textbf{twice}. Chapter~\ref{ch:blockchains} shows what that
coordination cost actually looks like; it is the dominant term.

The generalizable rule is that hybrid is cheap where the redundant work is
\emph{ephemeral, negotiated, and invisible}, and expensive where it is
\emph{persistent, consensus-fixed, and user-visible}. Before adopting hybrid by
reflex, check which of those two descriptions the system fits.

\section{Post-quantum TLS 1.3}
TLS 1.3 secures most web traffic, and it is the first place post-quantum cryptography has reached production. Two independent parts of the handshake must be upgraded: the \emph{key exchange} (confidentiality) and the \emph{authentication} (signatures on certificates).

\subsection{Hybrid key exchange}
As of this book's standards snapshot, TLS hybrid ECDHE--ML-KEM groups are under IETF standardization and are appearing in implementations. A group such as \texttt{X25519MLKEM768} combines an X25519 share with an ML-KEM contribution; the exact wire format, group naming, enablement, and TLS key-schedule binding must follow the applicable IETF specification and implementation release notes.

\begin{figure}[H]
\centering
\begin{tikzpicture}[font=\scriptsize]
  \node[font=\small\bfseries] at (1.4,3.4) {Client};
  \node[font=\small\bfseries] at (8.6,3.4) {Server};
  \draw[pqcgray!45, dashed] (1.4,3.1) -- (1.4,0.0);
  \draw[pqcgray!45, dashed] (8.6,3.1) -- (8.6,0.0);
  \draw[flowarrow] (1.4,2.5) -- node[above]{X25519 share $+$ ML-KEM key (encapsulation key)} (8.6,2.5);
  \draw[flowarrow] (8.6,1.7) -- node[above]{X25519 share $+$ ML-KEM ciphertext $+$ certificate} (1.4,1.7);
  \node[keynode, align=center] at (5,0.55) {shared secret $=\mathrm{KDF}\big(ss_{\text{X25519}} \,\|\, ss_{\text{ML-KEM}}\big)$};
\end{tikzpicture}
\caption{Illustrative hybrid key exchange in TLS 1.3. Both an elliptic-curve share and an ML-KEM contribution travel in the handshake; an analysed KDF construction binds both secrets and the protocol context. Post-quantum certificate authentication is a separate, draft/experimental deployment track. Specific group names and deployment status are version-dependent.}
\label{fig:tls-hybrid}
\end{figure}

\subsection{Post-quantum authentication and its size cost}
Post-quantum certificate authentication is still a separate standardization and
deployment track. ML-DSA-65 is one plausible illustrative choice for its balance of
speed and size, while SLH-DSA offers a more conservative assumption at larger sizes.
The practical obstacle is size: post-quantum public keys and signatures are far larger
than the elliptic-curve ones they replace, and a TLS handshake may carry several
certificates. The extra kilobytes cost round-trips and can collide with buffer and
fragmentation limits.

\begin{table}[H]
\centering
\begin{tabular}{lrr}
\toprule
Signature scheme & Public key (bytes) & Signature (bytes) \\
\midrule
Ed25519 (classical)   & 32   & 64 \\
ML-DSA-65 (FIPS 204)  & 1952 & 3309 \\
FN-DSA / Falcon-512 (FIPS 206 in development) & 897 & $\approx 666$ \\
SLH-DSA-128s (FIPS 205) & 32 & $\approx 7856$ \\
\bottomrule
\end{tabular}
\caption{Approximate sizes of a classical signature, published PQC standards, and the draft-dependent FN-DSA candidate. Falcon/FN-DSA is not a final FIPS at this standards snapshot.}
\label{tab:sig-sizes}
\end{table}

\subsection{What has actually shipped}
\label{sec:pq-adoption}

It is easy to read a chapter like this one as describing a future. For key
exchange, it does not. Hybrid post-quantum key agreement is on by default in
mainstream browsers and in OpenSSH~\cite{opensshpq}, and the share of
human-initiated HTTPS traffic reaching Cloudflare under a post-quantum key
agreement passed roughly one half during 2025~\cite{cloudflarepq2025}. Major
cloud key-management and transport services offer post-quantum KEMs as a
supported option. A reader who has browsed the web or opened an SSH session
recently has, with high probability, already used ML-KEM.

Authentication is the opposite story, and the contrast is the useful part.
Signature migration has barely begun: post-quantum certificates remain a
draft and experimental track, public certificate authorities are only starting
to issue them, and the chains of trust behind code signing, firmware, and secure
boot turn over on a timescale of years. The asymmetry has a simple cause.

\begin{itemize}
  \item Key exchange is \textbf{ephemeral and two-party}. Both endpoints
        negotiate afresh on every connection, so an upgrade propagates as fast as
        software updates do, and nothing needs to be re-issued.
  \item Signatures are \textbf{long-lived and multi-party}. A root certificate,
        a signed firmware image, or an on-chain public key commits a verifier who
        may not be upgraded --- or may not even exist --- yet.
\end{itemize}

That distinction recurs for the rest of this book. It is why ``harvest now,
decrypt later'' is the urgent framing for confidentiality but the wrong framing
for signatures, and it is why the hardest case in the next chapter is a
signature case.

Implementation and deployment status changes rapidly. Before making a production decision, consult the current IETF document and the release notes of the selected browser, TLS library, or service for its supported groups, defaults, and authentication behaviour.

\section{Verifying post-quantum signatures on-chain}
\label{sec:onchain}
A very different deployment target is the blockchain, where a smart contract may need to \emph{verify} a post-quantum signature -- for a quantum-resistant wallet, or a cross-chain bridge. Here the constraint is not bandwidth but \emph{gas}: every operation the Ethereum Virtual Machine (EVM) executes costs money, and a lattice verifier does a great deal of modular arithmetic.

Direct verification of a post-quantum signature can be expensive in an EVM-like environment because it requires substantial hashing and modular arithmetic. The cost depends on the scheme, parameter set, contract implementation, chain, hard-fork gas schedule, and available precompiles. Claims about a dominant operation or a gas figure require a reproducible benchmark with those details.

When even that is too expensive, an alternative is to avoid running the verifier on-chain at all.

\begin{itemize}
  \item \textbf{Direct verification}: run the full verifier in the contract. It is conceptually simple and minimizes added trust assumptions, but may be expensive.
  \item \textbf{ZK proof of verification}: prove off-chain that a signature verifies, then check a succinct proof on-chain. This moves work to the prover and introduces proof-system assumptions and costs.
  \item \textbf{Optimistic or fraud-proof designs}: accept a claim subject to a challenge procedure. Their cost and trust trade-offs depend on the exact challenge game and security model.
\end{itemize}

\begin{figure}[H]
\centering
\begin{tikzpicture}[node distance=5mm and 10mm, every node/.style={font=\scriptsize}]
  \node[flownode] (direct) {direct on-chain verify\\[-2pt]\scriptsize highest gas, simplest};
  \node[flownode, right=of direct] (zk) {ZK proof of verify\\[-2pt]\scriptsize cheap check, costly prove};
  \node[keynode, right=of zk] (nay) {optimistic / fraud proof\\[-2pt]\scriptsize needs challenge design};
  \draw[flowarrow] (direct) -- (zk);
  \draw[flowarrow] (zk) -- (nay);
\end{tikzpicture}
\caption{Three broad designs for on-chain post-quantum signature verification. They trade direct execution cost against prover work, challenge mechanisms, and security assumptions; there is no universal cost ordering without a specified benchmark and protocol.}
\label{fig:onchain-verify}
\end{figure}

\section{Sizes, performance, and implementation reality}
Choosing a scheme in practice is mostly a negotiation among sizes, speed, and assumption strength.

\begin{table}[H]
\centering
\begin{tabular}{llrr}
\toprule
Scheme & Role & Public key & Ciphertext / signature \\
\midrule
ML-KEM-768 (FIPS 203) & KEM & 1184 & 1088 \\
ML-DSA-65 (FIPS 204)  & signature & 1952 & 3309 \\
FN-DSA / Falcon-512 (FIPS 206 in development) & signature & 897 & $\approx 666$ \\
SLH-DSA-128s (FIPS 205) & signature & 32 & $\approx 7856$ \\
\bottomrule
\end{tabular}
\caption{Approximate sizes (in bytes) of published standards and the draft-dependent FN-DSA candidate. There is no universally best choice: SLH-DSA minimizes assumptions but has large signatures, while ML-DSA is a general-purpose published standard.}
\label{tab:scheme-sizes}
\end{table}

Two constraints cut across all of these. First, operations that handle secret keys,
ephemeral secrets, or secret-dependent sampling -- especially key generation, signing,
and decapsulation -- require side-channel review against the applicable threat model
(Chapter~\ref{ch:sidechannel}). Public-input verification has different confidentiality
requirements, though it still needs robustness and denial-of-service review. Second,
the sizes in Table~\ref{tab:scheme-sizes} are not academic -- they determine whether a
scheme fits in a TLS handshake, a certificate chain, or a smart-contract call.

\section{Pitfalls}

\paragraph{Pitfall 1: deploying post-quantum only, too early.}
Dropping the classical algorithm and trusting the post-quantum one alone forfeits the hedge that hybrid provides. Until the post-quantum schemes have had many more years of scrutiny, hybrid is the conservative default --- in the protocols where hybrid is cheap. Section~\ref{sec:hybrid-limits} gives the test for when it is not, and the mirror-image mistake, adopting hybrid where it doubles a permanent cost and forces users to migrate twice, is just as real.

\paragraph{Pitfall 2: ignoring handshake and message sizes.}
Post-quantum keys and signatures are kilobytes, not tens of bytes. Buffers, packet-size assumptions, and certificate-chain limits that were fine for elliptic curves can silently break; sizes must be designed for, not discovered in production.

\paragraph{Pitfall 3: assuming a reference implementation is deployable.}
Reference code optimizes for clarity, not for constant-time behaviour. Shipping it unmodified can reintroduce exactly the side channels of Chapter~\ref{ch:sidechannel}.

\paragraph{Pitfall 4: hardcoding one algorithm.}
The standards will keep changing. A system that cannot negotiate or swap algorithms will have to be re-engineered at the worst possible time.

\section{Summary}
The mathematics of the previous parts defeats Shor's algorithm, but that is where the abstract scheme ends and the engineering begins:

\begin{itemize}
  \item \textbf{hybrid} classical-plus-postquantum deployment, with \textbf{crypto-agility}, is the mainstream migration strategy --- where the redundant work is ephemeral, negotiated, and invisible, and not where it is permanent, consensus-fixed, and user-visible;
  \item \textbf{TLS 1.3} hybrid key exchange is under standardization and deployment; protocol and implementation versions must be checked at the time of use;
  \item post-quantum \textbf{key exchange has already shipped} at scale, while \textbf{signature} migration has barely started --- because key exchange is ephemeral and two-party, and signatures are long-lived and multi-party;
  \item verifying signatures \textbf{on-chain} has scheme- and environment-dependent costs, with direct, proof-based, and optimistic designs offering different trade-offs;
  \item scheme choice is a trade-off among size, speed, and assumption strength, with side-channel controls scoped to the secrets and threat model of each operation.
\end{itemize}

Post-quantum cryptography is no longer a paper exercise: the standards exist, and the migration is under way. Real security will arrive only when these schemes are deployed correctly, agilely, and side-channel-safely across the protocols and systems that the world already depends on. The next chapter turns to one of the hardest cases of all -- the blockchains, where the funds themselves are secured by exactly the signatures now at risk.
